% Section 2 -- Background: the syllable on the page, three encodings, noi lai. Budget: 0.75 page.
% Follows docs/DESIGN_DECISIONS.md sections 2, 3 (six-type taxonomy, region-neutral labels,
% old-style placement as the majority convention). Related work is Section 8; the full rule
% tables are Appendix B. Claim ledger: paper/claims.md rows L1-L12, N1, D13.
% Every Vietnamese example below is a row of data/attested_seed.tsv or an engine output
% (DESIGN_DECISIONS 3.6); all await native-speaker verification (the NV placeholder).
\section{Background}
\label{sec:background}

\paragraph{The Vietnamese syllable on the page.}
A written syllable is onset + rime + tone, the rime being an optional glide, a nucleus and an optional coda \citep{thompson-1965-vietnamese,doan-1977-nguam,nguyen-1997-vietnamese}; our parser recognizes 24 onsets (including the empty one), 18 nuclei, 11 codas (including none) and six tones, \vi{ngang} (unmarked), \vi{huyền}, \vi{sắc}, \vi{hỏi}, \vi{ngã} and \vi{nặng} \citep{kirby-2011-vietnamese,pham-2003-vietnamese-tone}. Structure and letters are not one-to-one: /k/ is written \vi{q} before the glide, \vi{k} before \vi{e ê i y} and \vi{c} elsewhere, and syllables ending in a stop take only \vi{sắc} and \vi{nặng} (Appendix~\ref{app:rules}). Two conventions place the tone mark; they differ only on the open rimes \vi{oa}, \vi{oe} and \vi{uy} after an onset other than \vi{qu}, where the older convention marks the first letter (\vi{hòa}, \vi{khỏe}) and the newer one the second (\vi{hoà}, \vi{khoẻ}): 69 of the 6,611 lower-case entries of the standard spelling list \citep{hunspell-vi}. Both are in everyday use; the older one is the majority in the Vietnamese test set we re-encode (51 of 53 affected tokens) and, as far as we can tell, in running text, the newer one the textbook convention (\placeholder{corpus count on two registers before Gate~1 fixes the majority and the direction of \hyp{4}}).

\paragraph{Three encodings of the same letters.}
A Vietnamese letter carries up to two diacritics, a quality mark (\vi{ă â ê ô ơ ư}) and a tone mark. Unicode has a precomposed code point for every combination, so the same text can be stored in NFC (one code point per letter), in NFD (base letter plus combining marks, with the tone mark \emph{before} the quality mark in \vi{ộ ậ ặ ệ}) or in a partially composed form inherited from Windows-1258 (quality mark precomposed, tone mark combining) that survives in legacy web text. All three are valid and occur in the wild. A tokenizer that normalizes to NFC erases the difference; one that passes its input through receives a different token sequence for the same word; some implementations \emph{corrupt} combining marks, deleting them. Neither Gemma tokenizer we audited normalizes (Section~\ref{sec:setup}).

\paragraph{Nói lái.}
\vi{Nói lái} is the Vietnamese spoonerism: two syllables exchange components so that a new, often humorous or hidden, phrase appears. Writing a syllable as $(O, R, T)$, a variant is a subset of $\{O, R, T\}$ exchanged between the two positions; the empty subset is the identity, the full one plain reversal, and the six others are the six types of the tradition \citep{le-ho-1990-thu-choi-chu,nguyen-van-hiep-noi-lai}: \var{1} swaps the rimes (\vi{mèo cái} $\rightarrow$ \vi{mài kéo}; \vi{trò chơi} `game' $\rightarrow$ \vi{trời cho} `heaven-sent'); \var{2} swaps onset and rime, so the syllables trade places and the tones stay (\vi{đầu tiên} `first' $\rightarrow$ \vi{tiền đâu} `where's the money?'); \var{3} swaps the tones (\vi{bí mật} `secret' $\rightarrow$ \vi{bị mất} `got lost'); \var{4} swaps rime and tone (\vi{bí mật} $\rightarrow$ \vi{bật mí} `reveal'); \var{5} swaps the onsets (\vi{giải pháp} `solution' $\rightarrow$ \vi{phải giáp} `must face') and \var{6} onsets and tones. A textbook illustration runs all six on \vi{thay đổi} `to change' (\vi{thôi đảy}, \vi{đôi thảy}, \vi{thảy đôi}, \vi{thổi đay}, \vi{đay thổi}, \vi{đảy thôi}), each of which our engine reproduces (\placeholder{confirm against \citet{le-ho-1990-thu-choi-chu}, not yet consulted directly}). Each variant is an involution, the eight form the group $(\mathbb{Z}/2)^3$, and reading an output in the other syllable order maps \var{1} to \var{6}, \var{2} to \var{3} and \var{4} to \var{5}; we therefore generate items for \var{1}--\var{4}, score \task{T1} in the named order and as an unordered pair, and keep \var{5} and \var{6} in the taxonomy, the \task{T2} gold and the \task{T3} twins (Section~\ref{sec:tasks}). When components coincide the variants collapse: \vi{thi đua} `to compete' $\rightarrow$ \vi{thua đi} `just lose' has equal tones, so \var{1} and \var{4} agree and \var{3} is the identity. Attested practice bends outputs toward real words: \vi{thầy giáo} `teacher' $\rightarrow$ \vi{tháo giày} `take off shoes' keeps the \vi{ay} of the real word where the rule gives \vi{giầy}, \vi{Vũ Như Cẩn} $\rightarrow$ \vi{vẫn như cũ} `still the same' is exact only under the Southern merger of \vi{hỏi} and \vi{ngã}, and three-syllable phrases swap a non-adjacent pair (\vi{chà đồ nhôm} `polish aluminium' $\rightarrow$ \vi{chôm đồ nhà} `steal from home'). The literature attaches regional labels to the variants inconsistently, so our codes are region-neutral and validators are region-tagged. Nói lái runs from folk sayings to the poetry of Hồ Xuân Hương, where it hides what a censor would not print; some attested examples have vulgar readings, which our release flags and keeps out of every human form and API prompt (see the Ethical Considerations section).\footnote{\placeholder{[NATIVE-CHECK] Every Vietnamese example in this paper is a row of the released attested seed table or an engine output on a dictionary pair; all forms and glosses await verification by the native validators.}}
